\documentclass[aspectratio=169,11pt]{beamer}

\usetheme{Madrid}
\usecolortheme{default}
\usefonttheme{professionalfonts}
\setbeamertemplate{navigation symbols}{}
\setbeamertemplate{footline}[frame number]

\usepackage{amsmath, amssymb, booktabs}

\title{Technology, Inequality, Market Structure, And Labor-Market Institutions}
\subtitle{Market-Level Incidence Of Technological Change}
\author{Special Topic 5: Technology, Innovation, and Labor -- Week 5}
\date{}

\begin{document}

\begin{frame}
  \titlepage
\end{frame}

\begin{frame}{Central question}
\begin{itemize}
    \item When technology raises productivity, who captures the gains?
    \item Do gains appear as wages, lower prices, rents, markups, or firm growth?
    \item Which workers move into better jobs, and which workers move into worse job-quality regimes?
    \item Week 5 discipline: market incidence is a labor question.
\end{itemize}
\end{frame}

\begin{frame}{From firm adoption to market incidence}
\small
\[
\Delta Y^{L}_{m}
=
\sum_f s_{fm}\Delta y^{L}_{fm}
+
\sum_f \Delta s_{fm}y^{L}_{fm}
+
\sum_j \Delta \omega_{jm}y^{L}_{jm}
\]
\begin{itemize}
    \item Within firms: productivity, wages, labor share, training, worker voice.
    \item Between firms: market share shifts toward different labor-share and wage-setting regimes.
    \item Between sectors: workers move across agriculture, manufacturing, formal services, informal services, and digital services.
    \item The same technology can raise output and worsen worker incidence.
\end{itemize}
\end{frame}

\begin{frame}{Rent capture, labor share, and superstar firms}
\begin{itemize}
    \item Technology can create rents through patents, data, software, organizational capital, and scale.
    \item Superstar-firm logic: market share shifts toward high-productivity, low-labor-share firms.
    \item Aggregate labor share can fall even without a representative firm becoming less labor friendly.
    \item Labor questions: which workers share rents, which workers lose outside options, and which jobs expand?
\end{itemize}
\end{frame}

\begin{frame}{Market power and institutions}
\small
\begin{tabular}{p{0.23\linewidth} p{0.37\linewidth} p{0.27\linewidth}}
\toprule
Margin & Technology channel & Labor incidence \\
\midrule
Employer concentration & growth of dominant adopters & weaker outside options \\
Intangible rents & data, patents, platforms, brands & uneven rent sharing \\
Platform rules & matching, ratings, scheduling & risk and control shift \\
Worker voice & unions, bargaining, transparency & incidence mediation \\
\bottomrule
\end{tabular}
\vspace{0.35cm}
\begin{itemize}
    \item Institutions do not change the engineering frontier, but they change who captures surplus.
\end{itemize}
\end{frame}

\begin{frame}{Historical technology waves}
\small
\begin{tabular}{p{0.20\linewidth} p{0.30\linewidth} p{0.35\linewidth}}
\toprule
Wave & Main labor margin & Persistence question \\
\midrule
Mechanization & manual effort and factory scale & did displaced tasks return elsewhere? \\
Electrification & workflow and organization & how long did complements take? \\
Computerization & routine cognitive tasks & did polarization persist? \\
Robots & embodied production tasks & did local losses offset firm gains? \\
AI & prediction, language, coding, control & new tasks or deeper reallocation? \\
\bottomrule
\end{tabular}
\vspace{0.3cm}
\begin{itemize}
    \item AI is new in capabilities, but it still must be read through task reallocation, diffusion, complements, and institutions.
\end{itemize}
\end{frame}

\begin{frame}{Structural change and job quality}
\begin{itemize}
    \item Agriculture to manufacturing to services is not only a productivity story.
    \item Mature economies face manufacturing-to-services transitions after automation and trade shocks.
    \item Employment reallocation is not the same as welfare-improving reallocation.
    \item Job-quality objects: wages, stability, voice, schedules, safety, formality, and mobility ladders.
    \item Formal services and informal services absorb workers in very different ways.
\end{itemize}
\end{frame}

\begin{frame}{Global and Global South evidence}
\begin{itemize}
    \item Developing-country technology incidence often runs through diffusion, absorptive capacity, informality, and management.
    \item Labor-absorbing services can be a development path or a low-productivity absorption margin.
    \item Outsourcing and offshoring move tasks across countries before all jobs are automated.
    \item AI may complement scarce skilled labor in some settings and displace digital service tasks in others.
    \item High-income automation evidence is not automatically externally valid.
\end{itemize}
\end{frame}

\begin{frame}{Frontier AI infrastructure questions}
\begin{itemize}
    \item Compute-intensive AI can shift electricity demand and data-center location.
    \item Data-center buildout affects construction, technical operations, local services, land use, and local fiscal capacity.
    \item Energy bottlenecks can change which firms and regions can adopt AI at scale.
    \item Outsourced digital services may expand through data work, localization, oversight, and AI-assisted knowledge services.
    \item The labor object is spatial incidence, not generic energy commentary.
\end{itemize}
\end{frame}

\begin{frame}{Methods and data}
\small
\begin{tabular}{p{0.33\linewidth} p{0.47\linewidth}}
\toprule
Design family & Labor object \\
\midrule
Market-share decomposition & within-firm versus between-firm labor-share change \\
Concentration accounting & product-market and labor-market power \\
Shift-share exposure & local incidence of technology shocks \\
Historical comparison & persistence, reversal, and complements \\
Structural-change decomposition & sectoral employment versus job quality \\
Linked firm-worker data & wage premia, training, exits, and mobility \\
Cross-country adoption surveys & diffusion and absorptive capacity \\
Energy and spatial data & AI infrastructure and local labor demand \\
\bottomrule
\end{tabular}
\end{frame}

\begin{frame}{Research Lab design}
\begin{itemize}
    \item Reproduce: build a synthetic industry-firm panel with market share, labor share, productivity, and technology intensity.
    \item Diagnose: separate direct technology measurement from inferred technology, concentration, wage setting, and institutions.
    \item Transfer: move from superstar firms to structural change, Global South diffusion, outsourcing, and AI infrastructure.
    \item Output: compact tables on labor-share decomposition, incidence diagnosis, and global/frontier transfer scenarios.
    \item Goal: practice market-level labor incidence without confidential firm microdata.
\end{itemize}
\end{frame}

\begin{frame}{Bridge to Week 6}
\begin{itemize}
    \item Week 5 aggregates technology from firms and workers into market-level incidence.
    \item Week 6 turns the whole course into research designs.
    \item A strong project names the technology margin, exposure measure, labor outcome, mediator, and counterfactual.
\end{itemize}
\end{frame}

\begin{frame}{Takeaways}
\begin{itemize}
    \item Technology changes who captures rents, not only how much output is produced.
    \item Labor share, concentration, firm rents, and wage setting must be studied together.
    \item Historical waves help discipline AI claims without flattening what is new.
    \item Structural change should be evaluated by job quality and welfare, not employment movement alone.
    \item Global evidence changes the interpretation of technology shocks and opens infrastructure and outsourcing frontiers.
\end{itemize}
\end{frame}

\end{document}
